% Part: incompleteness
% Chapter: introduction
% Section: definitions

\documentclass[../../../include/open-logic-section]{subfiles}

\begin{document}

\olfileid{inc}{int}{def}

\olsection{تعريفونه}

د رياضياتو د صوري کولو، او د داسې صوري کولو د سازګارۍ او بشپړتيا د
ښودلو د هېلبرټ پروژې د عملي کولو دپاره لومړنی کار دا و چې يوه ژبه،
منطقي چوکاټ او د بديهي اصولو نظام وټاکل شي۔ د خپلې موخې دپاره به فرض
کړو چې رياضيات د لومړۍ درجې په يوه ژبه کښې صوري کېدے شي؛ يعنې د
!!{constant}s، !!{function}s او !!{predicate}s داسې يو سټ شته چې د
لومړۍ درجې منطق له تړونکو او کميت‌ټاکونکو سره يوځاے موږ ته د رياضياتو
د ادعاوو د بيان وس راکوي۔ زياتره خلک مني چې داسې ژبه شته: د سټ تيورۍ
ژبه، چې~$\in$ پکې يوازينۍ غيرمنطقي نښه ده۔ دا چې دومره ساده ژبه دومره
بيانوونکې ده، البته په لومړي نظر ډېره نامعقوله ادعا ښکاري، او د دې د
ثابتولو دپاره ډېر کار وشو چې نږدې ټول رياضيات په همدې ډېرې محدودې
اصطلاح‌زېرمه کښې بيانېدے شي۔ د ساده‌والي دپاره به اوس خپل بحث حساب
ته محدود کړو، يعنې د رياضياتو هغې برخې ته چې يوازې له طبيعي
عددونو~$\Nat$ سره کار لري۔ د حساب د حقايقو د بيان طبيعي ژبه~$\Lang
L_A$ ده۔ $\Lang L_A$ يو دوه‌ځاييز !!{predicate}~$<$، يو
!!{constant}~$\Obj 0$، يوه يوځاييزه !!{function}~$\prime$، او دوه
دوه‌ځاييزې !!{function}s~$+$ او~$\times$ لري۔

\begin{defn}
د !!{sentence}s سټ~$\Gamma$ يوه \emph{تيوري} ده که د لازمې پايلې له
مخې تړلې وي، يعنې که
$\Gamma = \Setabs{!A}{\Gamma \Entails !A}$ وي۔
\end{defn}

د تيوريو د مشخصولو دوه اسانه لارې شته۔ يوه دا ده چې تيوري په يوه
!!{structure} کښې د رښتينو !!{sentence}s سټ وټاکو۔ د بېلګې په توګه،
د~$\Lang L_A$ هغه !!{structure} وګورئ چې !!{domain} ئې~$\Nat$ وي او
ټولې غيرمنطقي نښې پکې هغسې تفسير شوې وي لکه تمه چې ترې کېږي۔

\begin{defn}\ollabel{def:standard-model}
د حساب \emph{معياري مدل} هغه !!{structure}~$\Struct N$ دے چې داسې
تعريفېږي:
\begin{enumerate}
\item $\Domain N = \Nat$
\item $\Assign{\Obj 0}{N} = 0$
\item $\Assign{\Obj \prime}{N}(n) = n + 1$ د ټولو $n \in \Nat$ دپاره
\item $\Assign{\Obj +}{N}(n, m) = n + m$ د ټولو $n, m \in \Nat$ دپاره
\item $\Assign{\Obj \times}{N}(n, m) = n\cdot m$ د ټولو $n, m \in \Nat$ دپاره
\item $\Assign{\Obj <}{N} = \Setabs{\tuple{n, m}}{n \in \Nat, m \in
  \Nat, n < m}$
\end{enumerate}
\end{defn}

د~`$\times$' او~`$\cdot$' تر منځ توپير ته پام وکړئ: $\times$ د حساب
د ژبې يوه نښه ده۔ البته دا مو داسې ټاکلې چې ضرب راياد کړي، خو
$\times$ د ضرب عمل نۀ دے، بلکې يوه دوه‌ځاييزه تابع‌نښه ده (په رسمي
ډول~$\Obj f^2_1$)۔ د دې پر خلاف، $\cdot$ پخپله د ضرب عادي تابع
\emph{ده}۔ کله چې د~$n \cdot m$ په شان څه وينئ، موخه مو د~$n$ او~$m$
عددونو حاصل‌ضرب وي؛ کله چې د~$x \times y$ په شان څه وينئ، د حساب د
ژبې د يوې اصطلاح خبره کوو۔ په معياري مدل کښې تابع‌نښه~$\times$ پر
طبيعي عددونو د تابع~$\cdot$ په توګه تفسيرېږي۔ د جمع دپاره~$+$ هم د
حساب د ژبې د تابع‌نښې او هم پر طبيعي عددونو د جمع د تابع په توګه
کاروو۔ دلته بايد له چاپېرياله معلومه کړئ چې کومه معنا موخه ده۔

\begin{defn}
د \emph{رښتيني حساب} تيوري د هغو !!{sentence}s سټ دے چې د حساب په
معياري مدل کښې پوره کېږي، يعنې
\[
\Th{TA} = \Setabs{!A}{\Sat{N}{!A}}.
\]
\end{defn}

$\Th{TA}$ يوه تيوري ده، ځکه هر کله چې~$\Th{TA} \Entails !A$ وي،
$!A$ په هر هغه !!{structure} کښې پوره کېږي چې~$\Th{TA}$ پوره کوي۔
څرنګه چې~$\Sat{N}{\Th{TA}}$، نو~$\Sat{N}{!A}$ لرو، او له دې امله
$!A \in \Th{TA}$۔

د تيورۍ~$\Gamma$ د مشخصولو بله لار دا ده چې هغه د هغو !!{sentence}s سټ وټاکو
چې د جملو له يوه سټ~$\Gamma_0$ څخه لازمېږي۔ په دې حالت کښې~$\Gamma$
د لازمې پايلې له مخې د~$\Gamma_0$ «تړون» دے۔ د تيورۍ دا ډول مشخصول
يوازې هغه وخت په زړه پورې دي چې~$\Gamma_0$ په څرګند ډول مشخص شوے وي؛
مثلاً د~$\Gamma_0$ !!{element}s لست شوي وي۔ لږ تر لږه~$\Gamma_0$
بايد د پرېکړې وړ وي؛ يعنې داسې محاسبه کېدونکې ازموينه بايد وي چې يوه
!!a{sentence} د~$\Gamma_0$ غړې ده که نۀ۔ د~$\Gamma_0$ !!{sentence}s
د~$\Gamma$ دپاره \emph{بديهي اصول} بولو، او وايو چې~$\Gamma$ د
$\Gamma_0$ په وسيله \emph{بديهي شوې} ده۔

\begin{defn}
تيوري~$\Gamma$ هله او يوازې هله د~$\Gamma_0$ په وسيله بديهي شوې ده
چې
\[
\Gamma = \Setabs{!A}{\Gamma_0 \Entails !A}
\]
\end{defn}

\begin{defn}
هغه تيوري~$\Th{Q}$ چې په لاندې جملو بديهي شوې ده د «رابنسن
$\Th{Q}$» په نوم پېژندل کېږي او د حساب يوه ډېره ساده تيوري ده۔
\begin{align*}
& \lforall[x][\lforall[y][(\eq[x'][y'] \lif \eq[x][y])]] \tag{$!Q_1$}\\
& \lforall[x][\eq/[\Obj 0][x']] \tag{$!Q_2$}\\
& \lforall[x][(\eq[x][\Obj 0] \lor \lexists[y][\eq[x][y']])] \tag{$!Q_3$}\\
& \lforall[x][\eq[(x + \Obj 0)][x]] \tag{$!Q_4$}\\
& \lforall[x][\lforall[y][\eq[(x + y')][(x + y)']]] \tag{$!Q_5$}\\
& \lforall[x][\eq[(x \times \Obj 0)][\Obj 0]] \tag{$!Q_6$}\\
& \lforall[x][\lforall[y][\eq[(x \times y')][((x \times y) + x)]]] \tag{$!Q_7$}\\
& \lforall[x][\lforall[y][(x < y \liff \lexists[z][\eq[(z' + x)][y]])]] \tag{$!Q_8$}
\end{align*}
د !!{sentence}s سټ~$\{!Q_1, \dots, !Q_8\}$ د~$\Th{Q}$ بديهي اصول
دي، نو~$\Th{Q}$ له هغو څخه لازمېدونکي ټول !!{sentence}s رانغاړي:
\[
\Th{Q} = \Setabs{!A}{\{!Q_1, \dots, !Q_8\} \Entails !A}.
\]
\end{defn}

\begin{defn}
فرض کړئ~$!A(x)$ په~$\Lang L_A$ کښې داسې !!a{formula} ده چې آزاد
متغيرونه~$x$ او~$y_1$، \dots،~$y_n$ لري۔ بيا د لاندې بڼې هره
!!{sentence}
\[
\lforall[y_1][\dots\lforall[y_n][((!A(\Obj 0) \land \lforall[x][(!A(x)
\lif !A(x'))]) \lif \lforall[x][!A(x)])]]
\]
د \emph{استقرا شېما} يوه بېلګه ده۔

\emph{د پيانو حساب}~$\Th{PA}$ هغه تيوري ده چې د~$\Th{Q}$ په بديهي
اصولو او د استقرا شېما په ټولو بېلګو بديهي شوې ده۔
\end{defn}

\begin{explain}
د استقرا شېما هره بېلګه په~$\Struct{N}$ کښې رښتونې ده۔ دا خبره تر
ټولو اسانه هغه وخت ليدل کېږي چې !!{formula}~$!A$ يوازې يو آزاد
!!{variable}~$x$ ولري۔ بيا~$!A(x)$ په~$\Struct{N}$ کښې د~$\Nat$ يو
فرعي سټ~$X_{!A}$ تعريفوي۔ $X_{!A}$ د ټولو هغو~$n \in \Nat$ سټ دے
چې~$s(x) = n$ وي او~$\Sat{N}{!A(x)}[s]$۔ د استقرا شېما اړونده بېلګه
دا ده:
\[
((!A(\Obj 0) \land \lforall[x][(!A(x) \lif !A(x'))]) \lif
  \lforall[x][!A(x)]).
\]
که مقدمه ئې په~$\Struct{N}$ کښې رښتونې وي، نو~$0 \in X_{!A}$، او
هر کله چې~$n \in X_{!A}$ وي، $n+1$ هم پکې وي۔ له~$0 \in X_{!A}$
څخه~$1 \in X_{!A}$ ترلاسه کوو۔ له~$1 \in X_{!A}$ څخه~$2 \in X_{!A}$
ترلاسه کوو، او همداسې وړاندې۔ نو د هر~$n \in \Nat$ دپاره
$n \in X_{!A}$۔ خو دا په دې معنا ده چې~$\lforall[x][!A(x)]$ په
$\Struct{N}$ کښې پوره کېږي۔
\end{explain}

$\Th{Q}$ او~$\Th{PA}$ دواړه بديهي شوې تيورۍ دي۔ لويه پوښتنه دا ده
چې څومره پياوړې دي؟ مثلاً آيا~$\Th{PA}$ د~$\Nat$ په باب ټول هغه
حقايق ثابتولے شي چې په~$\Lang L_A$ کښې بيانېدے شي؟ په کره توګه، آيا
د~$\Th{PA}$ بديهي اصول ټولې هغه پوښتنې حل کوي چې په~$\Lang L_A$ کښې
جوړېدے شي؟

د دې پوښتنې بله بڼه دا ده: آيا~$\Th{PA} = \Th{TA}$؟ $\Th{TA}$ په
څرګند ډول د~$\Nat$ په باب ټول حقايق ثابتوي (يعنې ټول رانغاړي)، او
ټولې هغه پوښتنې حل کوي چې په~$\Lang L_A$ کښې جوړېدے شي؛ ځکه که~$!A$
په~$\Lang L_A$ کښې !!a{sentence} وي، نو يا~$\Sat{N}{!A}$ يا
$\Sat{N}{\lnot !A}$، او له دې امله يا~$\Th{TA} \Entails !A$ يا
$\Th{TA} \Entails \lnot !A$۔ داسې تيوري \emph{!!{complete}} بولو۔

\begin{defn}
تيوري~$\Gamma$ هله او يوازې هله \emph{!!{complete}} ده چې د هغې د
ژبې د هرې !!{sentence}~$!A$ دپاره يا~$\Gamma \Entails !A$ يا
$\Gamma \Entails \lnot !A$ وي۔
\end{defn}

\begin{explain}
د بشپړتيا د قضيې له مخې~$\Gamma \Entails !A$ هله او يوازې هله چې
$\Gamma \Proves !A$، نو~$\Gamma$ هله او يوازې هله بشپړه ده چې د
هغې د ژبې د هرې !!{sentence}~$!A$ دپاره يا~$\Gamma \Proves !A$ يا
$\Gamma \Proves \lnot !A$ وي۔
\end{explain}

بله پوښتنه چې راپورته کېږي دا ده: آيا داسې محاسبوي کړنلاره شته چې
وآزمويو يوه !!a{sentence} په~$\Th{TA}$، په~$\Th{PA}$، يا ان يوازې
په~$\Th{Q}$ کښې ده که نۀ؟ دا پوښتنه د دې په تعريفولو لا کره کولے شو
چې يو سټ (مثلاً د !!{sentence}s سټ) کله !!{decidable} وي۔

\begin{defn}
سټ~$X$ هله او يوازې هله \emph{!!{decidable}} دے چې داسې محاسبوي
کړنلاره وي چې پر ننوت~$x$، که~$x \in X$ وي~$1$ او که نۀ وي~$0$
بېرته ورکوي۔
\end{defn}

نو زموږ پوښتنه داسې شوه: آيا~$\Th{TA}$ ($\Th{PA}$، $\Th{Q}$)
!!{decidable} ده؟

د دې ټولو پوښتنو ځواب به «نه» وي۔ له دې تيوريو څخه يوه هم د پرېکړې
وړ نۀ ده۔ خو دا پديده يوازې په همدې ځانګړو تيوريو پورې محدوده نۀ ده۔
په حقيقت کښې \emph{هره} هغه تيوري چې ځينې شرطونه پوره کوي له همدغو
پايلو سره مخ کېږي۔ له دغو شرطونو څخه يو، چې~$\Th{Q}$ او~$\Th{PA}$
ئې پوره کوي، دا دے چې د بديهي اصولو د !!a{decidable} سټ په وسيله
بديهي شوې وي۔

\begin{defn}
يوه تيوري \emph{!!{axiomatizable}} ده که د بديهي اصولو د
!!a{decidable} سټ په وسيله بديهي شوې وي۔
\end{defn}

\begin{ex}
هره تيوري چې د !!{sentence}s د يوه متناهي سټ په وسيله بديهي شوې وي
!!{axiomatizable} ده، ځکه هر متناهي سټ !!{decidable} دے۔ نو مثلاً
$\Th{Q}$، !!{axiomatizable} ده۔

د~$\Th{PA}$ په شان په شېمايي ډول بديهي شوې تيورۍ هم
!!{axiomatizable} دي۔ د دې د ازمويلو دپاره چې~$!B$ د~$\Th{PA}$ په
بديهي اصولو کښې شامله ده که نۀ، يعنې د تابع~$\Char{X}$ د محاسبې
دپاره، چې که~$!B$ د~$\Th{PA}$ يو بديهي اصل وي
$\Char{X}(!B) = 1$ او که نۀ وي~$=0$، داسې کولے شو: لومړے وګورئ
$!B$ د~$\Th{Q}$ له بديهي اصولو څخه ده که نۀ۔ که وي، ځواب «هو» دے
او~$\Char{X}(!B) = 1$۔ که نۀ وي، وازمويئ چې د استقرا شېما بېلګه ده
که نۀ۔ دا په منظمه توګه کېدے شي؛ په دې حالت کښې ښايي تر ټولو اسانه
دا وي چې داسې ئې ووينو: د استقرا شېما هره بېلګه په يو شمېر کلي
کميت‌ټاکونکو پيلېږي، او بيا داسې فرعي-!!{formula} راځي چې يو شرطي
دے۔ د هغه شرطي پايله~$\lforall[x][!A(x, y_1, \dots, y_n)]$ ده، چې
$x$ او~$y_1$، \dots،~$y_n$ د~$!A$ ټول آزاد متغيرونه دي او د~$!B$
په پيل کښې کميت‌ټاکونکي متغيرونه~$y_1$، \dots،~$y_n$ تړي۔ کله چې دا
$!A$ راوباسو او وګورو چې آزاد متغيرونه ئې د لومړنيو کلي
کميت‌ټاکونکو او~$\lforall[x]$ له تړلو متغيرونو سره برابر دي، بيا
ګورو چې د شرطي مقدمه له لاندې سره برابره ده که نۀ:
\[
!A(\Obj 0, y_1, \dots, y_n) \land \lforall[x][(!A(x, y_1, \dots, y_n)
\lif !A(x', y_1, \dots, y_n))]
\]
بيا هم، که برابره وي، $!B$ د استقرا شېما بېلګه ده؛ او که برابره نۀ
وي، $!B$ د استقرا شېما بېلګه نۀ ده۔
\end{ex}

د دې پوښتنې، او د دې لا عامې پوښتنې د ځوابولو دپاره چې کومې تيورۍ
بشپړې يا د پرېکړې وړ دي، لاندې تعريف هم ګټور دے۔ په ياد ولرئ چې سټ~$X$
هله او يوازې هله !!{enumerable} دے چې تش وي يا داسې !!a{surjective}
تابع~$f \colon \Nat \to X$ موجوده وي۔ داسې تابع د~$X$ شمېره بلل کېږي۔

\begin{defn}
سټ~$X$ هله او يوازې هله \emph{!!{computably enumerable}} (په لنډ ډول
!!{c.e.}) بلل کېږي چې تش وي يا محاسبه کېدونکې شمېره ولري۔
\end{defn}

د !!{axiomatizability} تر څنګ، پر هغو تيوريو بل شرط چې د ناتکميلۍ
قضيې پرې پلي کېږي دا دے چې د محاسبه کېدونکو تابعو او
!!{decidable} اړيکو په باب بنسټيز حقايق ثابتولو ته کافي پياوړې وي۔
له «بنسټيزو حقايقو» څخه موخه هغه !!{sentence}s دي چې د هرې تابعې د
ارګومېنټونو دپاره ئې ارزښتونه بيانوي۔ او له «کافي پياوړې» څخه موخه
دا ده چې اړوندې تيورۍ دا جملې په خپلو قضيو کښې شمېري۔ د بېلګې په
توګه يوه نمونه‌يي محاسبه کېدونکې تابع، يعنې جمع، وګورئ۔ د~$2$ او~$3$
ارګومېنټونو دپاره د~$+$ ارزښت~$5$ دے، يعنې~$2+3 = 5$۔ د حساب په
ژبه کښې يوه جمله چې وايي د~$2$ او~$3$ ارګومېنټونو دپاره د~$+$ ارزښت
$5$ دے، دا ده: $(\num{2} + \num{3}) = \num{5}$۔ او مثلاً~$\Th{Q}$
دا جمله ثابتوي۔ په عام ډول غواړو د هرې محاسبه کېدونکې تابع
$f(x_1, x_2)$ دپاره په~$\Lang{L_A}$ کښې داسې !!a{formula}
$!A_f(x_1, x_2, y)$ وي چې هر کله~$f(n_1, n_2) = m$،
$\Th{Q} \Proves !A_f(\num{n_1}, \num{n_2}, \num{m})$ وي۔ په دې ډول
$\Th{Q}$ ثابتوي چې د~$n_1$، $n_2$ ارګومېنټونو دپاره د~$f$ ارزښت~$m$
دے۔ په حقيقت کښې لږ څه زيات غواړو، يعنې دا هم ثابته کړي چې د~$n_1$،
$n_2$ ارګومېنټونو دپاره بل هېڅ عدد د~$f$ ارزښت نۀ دے۔ د
!!{decidable} اړيکو دپاره هم همداسې ده۔ لاندې دوه تعريفونه دا خبره
کره کوي۔

\begin{defn}
!!{formula}~$!A(x_1, \dots, x_k, y)$ په~$\Gamma$ کښې تابع
$f\colon \Nat^k \to \Nat$ هله او يوازې هله \emph{!!{represents}}
کوي چې هر کله~$f(n_1, \dots, n_k) = m$ وي، نو
\begin{enumerate}
\item $\Gamma \Proves !A(\num{n_1}, \dots, \num{n_k}, \num{m})$، او
\item $\Gamma \Proves \lforall[y](!A(\num{n_1}, \dots, \num{n_k},
y) \lif y = \num{m})$۔
\end{enumerate}
\end{defn}

\begin{defn}
!!{formula}~$!A(x_1, \dots, x_k)$ اړيکه~$R \subseteq \Nat^k$ هله او
يوازې هله \emph{!!{represents}} کوي چې
\begin{enumerate}
\item هر کله~$R(n_1, \dots, n_k)$ وي،
$\Gamma \Proves !A(\num{n_1}, \dots, \num{n_k})$، او
\item هر کله~$R(n_1, \dots, n_k)$ نۀ وي،
$\Gamma \Proves \lnot !A(\num{n_1}, \dots, \num{n_k})$۔
\end{enumerate}
\end{defn}

يوه تيوري د ناتکميلۍ د قضيو د پلي کېدو دپاره «کافي پياوړې» ده که
ټولې محاسبه کېدونکې تابعې او ټولې !!{decidable} اړيکې !!{represents}
کوي۔ $\Th{Q}$ او د هغې غځونې دا شرط پوره کوي، خو د دې ثابتول به لږ
وخت واخلي؛ دا د هغو خبرو په باب يو غيرابتذالي حقيقت دے چې~$\Th{Q}$
ئې ثابتولے شي، او ښودل ئې ځکه ستونزمن دي چې~$\Th{Q}$ يوازې څو بديهي
اصول لري چې بايد دا ټول حقايق ترې ثابت کړو۔ خو~$\Th{Q}$ ډېره کمزورې
تيوري ده۔ نو که څه هم دا ثابتول ستونزمن دي چې~$\Th{Q}$ ټولې محاسبه
کېدونکې تابعې تمثيلوي، زياتره په زړه پورې تيورۍ تر~$\Th{Q}$ پياوړې
دي؛ يعنې تر~$\Th{Q}$ زيات څه ثابتوي۔ او هر څه چې~$\Th{Q}$ ثابتوي،
هره پياوړې تيوري ئې هم ثابتوي؛ څرنګه چې~$\Th{Q}$ ټولې محاسبه کېدونکې
تابعې تمثيلوي، هره پياوړې تيوري ئې هم تمثيلوي۔ په دې معنا ډېرې په
زړه پورې تيورۍ د ناتکميلۍ د قضيو دا شرط پوره کوي۔ نو زموږ سخت کار به
ګټه وکړي، ځکه ښيي چې د ناتکميلۍ قضيې پر ډېرو تيوريو پلي کېږي۔ په
يقيني ډول هره هغه تيوري چې موخه ئې د «ټولو رياضياتو» صوري کول وي،
بايد هر څه ثابت کړي چې~$\Th{Q}$ ئې ثابتوي، ځکه لږ تر لږه بايد د
لومړنيو محاسبو پايلې رانغاړلے شي۔ نو هره هغه تيوري چې د «ټولو
رياضياتو» د تيورۍ نومانده وي، د ناتکميلۍ قضيې به پرې پلي کېږي۔


\end{document}
